\documentclass[10pt]{article}

% COMPILAR CON XeLaTeX o LuaLaTeX
\usepackage{fontspec}
\setmainfont{Fira Sans Light}
\setsansfont{Fira Sans Light}
\setmonofont{Fira Code}

% IDIOMA
\usepackage[spanish]{babel}

% MÁRGENES
\usepackage{geometry}
\geometry{a4paper, margin=2.2cm}

% COLORES
\usepackage[dvipsnames]{xcolor}

% MATEMÁTICAS
\usepackage{amsmath, amssymb, amsthm}
\usepackage{mathtools}

% TEOREMAS
\newtheorem{definition}{Definición}[section]
\newtheorem{theorem}{Teorema}[section]
\newtheorem{lemma}{Lema}[section]
\newtheorem{proposition}{Proposición}[section]
\newtheorem{corollary}{Corolario}[section]

% GRÁFICOS
\usepackage{tikz}
\usepackage{pgfplots}
\pgfplotsset{compat=1.18}

% LISTADOS
\usepackage{enumitem}

% CÓDIGO
\usepackage{listings}
\lstdefinestyle{pythonEstilo}{
    language=Python,
    basicstyle=\ttfamily\small,
    keywordstyle=\color{RoyalBlue}\bfseries,
    commentstyle=\color{green!50!black},
    stringstyle=\color{BrickRed},
    numbers=none,
    breaklines=true,
    frame=single
}

% TABLAS
\usepackage{booktabs}

% OTROS
\setlength{\parindent}{0pt}
\setlength{\parskip}{0.4em}


%---------------------------------------------
% DOCUMENTO
%---------------------------------------------
\begin{document}

\begin{center}
{\Large \bfseries Práctica 1 — Lectura y adquisición de datos de un receptor GPS}\[0.3em]
{\normalsize Sistemas de Posicionamiento por Satélite (SISGPS) — Facultad de Ingeniería — 2025}\[0.2em]
{\normalsize \textbf{Grupo:} \underline{\hspace{6cm}} \hfill \textbf{Fecha:} \today}
\end{center}

\vspace{0.5em}
\tableofcontents
\newpage


\section{Introducción}
En esta práctica se implementa una aplicación capaz de leer, a través de un puerto serie, tramas de posicionamiento según el estándar NMEA 0183 emitidas por un receptor GPS a 1~Hz. De dichas tramas se selecciona la sentencia \texttt{GGA} (\textit{Global Positioning System Fix Data}), que contiene la posición estimada por el receptor en coordenadas geográficas (latitud y longitud) junto con información de calidad del fix.

A partir de esas coordenadas geográficas se realiza la transformación a coordenadas UTM (\textit{Universal Transverse Mercator}), obteniendo la proyección plana en coordenadas cartesianas \((X, Y)\) en metros.

\section{Objetivo y entregables}
\textbf{Objetivo:} leer tramas \texttt{GGA} desde un receptor GPS por puerto serie y convertir la posición a UTM mediante las ecuaciones de Gauss--Krüger (Transversa de Mercator) indicadas en el enunciado.

\textbf{Entregables:}
\begin{itemize}[leftmargin=1.2em]
\item Código fuente de la aplicación (archivo \texttt{practica1.py}).
\item Documento PDF con la especificación del trabajo desarrollado (este informe).
\end{itemize}

\section{Formato de la trama GGA (NMEA 0183)}
El receptor emite tramas como:
\begin{center}
\texttt{\$GPGGA,123519,4807.038,N,01131.000,E,1,08,0.9,545.4,M,46.9,M,,*47}
\end{center}

En la Tabla~\ref{tab:gga} se resume el significado de los campos empleados por la aplicación.

\begin{table}[h]
\centering
\renewcommand{\arraystretch}{1.15}
\begin{tabular}{@{}ll@{}}
\toprule
\textbf{Campo} & \textbf{Descripción} \\
\midrule
\texttt{GGA} & Tipo de trama: \textit{Global Positioning System Fix Data} \\
\texttt{hhmmss.ss} & Hora UTC del fix \\
\texttt{lat, N/S} & Latitud en formato \texttt{ddmm.mmmm} y hemisferio \\
\texttt{lon, E/W} & Longitud en formato \texttt{dddmm.mmmm} y hemisferio \\
\texttt{fix} & Calidad del fix (0 = inválido, 1 = GPS, 2 = DGPS, \dots) \\
\texttt{sat} & Número de satélites seguidos \\
\texttt{HDOP} & Dilución horizontal de precisión \\
\texttt{alt} & Altitud sobre el nivel medio del mar (m) \\
\texttt{*CS} & Checksum XOR en hexadecimal (entre \texttt{\$} y \texttt{*}) \\
\bottomrule
\end{tabular}
\caption{Campos relevantes de la sentencia GGA.}
\label{tab:gga}
\end{table}

\subsection{Conversión de coordenadas NMEA a grados decimales}
La latitud y longitud se reciben en formato NMEA:
\begin{itemize}[leftmargin=1.2em]
\item Latitud: \texttt{ddmm.mmmm}
\item Longitud: \texttt{dddmm.mmmm}
\end{itemize}

La conversión a grados decimales se realiza como:
\[
\text{grados decimales} = \text{grados} + \frac{\text{minutos}}{60}
\]
Aplicando el signo según hemisferio: Sur (S) negativo y Oeste (W) negativo.

\section{Modelo elipsoidal y parámetros}
La transformación a UTM depende del elipsoide (datum). En la práctica se consideraron los dos elipsoides indicados en el enunciado: WGS84 y Europeo 1950 (ED50).

\begin{center}
\renewcommand{\arraystretch}{1.2}
\begin{tabular}{@{}lrrrr@{}}
\toprule
\textbf{Elipsoide} & $a$ (m) & $f$ & $e'^2$ & $c$ (m) \\
\midrule
Europeo 1950 (ED50) & 6378388 & $1/297$ & 0.00672267 & 6399936.608 \\
WGS84 & 6378137 & $1/298.25722$ & 0.00669437999013 & 6399593.626 \\
\bottomrule
\end{tabular}
\end{center}

Donde:
\begin{itemize}[leftmargin=1.2em]
\item $a$ es el semieje mayor.
\item $b$ es el semieje menor (se obtiene como $b=a(1-f)$).
\item $f$ es el achatamiento $f = \dfrac{a-b}{a}$.
\item $e'^2$ es el cuadrado de la segunda excentricidad $e'^2 = \dfrac{a^2-b^2}{b^2}$.
\item $c$ es el radio polar de curvatura $c = \dfrac{a^2}{b}$.
\end{itemize}

\section{Transformación a UTM (Gauss--Krüger / Transversa de Mercator)}
Sea $\varphi$ la latitud en radianes y $\lambda$ la longitud en radianes. El huso UTM se selecciona entre 1 y 60; para España se utiliza el \textbf{huso 30}.

\subsection{Meridiano central del huso}
\[
\lambda_{0g} = \textit{Huso}\cdot 6 - 183^\circ
\qquad
\lambda_0 = \lambda_{0g}\cdot \frac{\pi}{180}
\]

\subsection{Ecuaciones usadas}
En la implementación se emplean las expresiones del enunciado (forma estándar de TM/UTM). Se define:
\[
N = \frac{a}{\sqrt{1-e^2\sin^2\varphi}},\quad
T = \tan^2\varphi,\quad
C = e'^2\cos^2\varphi,\quad
A = \cos\varphi\,(\lambda-\lambda_0)
\]

El arco meridiano:
\[
M = a\Big[\alpha_0\varphi - \alpha_2\sin(2\varphi) + \alpha_4\sin(4\varphi) - \alpha_6\sin(6\varphi)\Big]
\]
con coeficientes:
\[
\alpha_0=1-\frac{e^2}{4}-\frac{3e^4}{64}-\frac{5e^6}{256},\quad
\alpha_2=\frac{3e^2}{8}+\frac{3e^4}{32}+\frac{45e^6}{1024},
\]
\[
\alpha_4=\frac{15e^4}{256}+\frac{45e^6}{1024},\quad
\alpha_6=\frac{35e^6}{3072}.
\]

Finalmente, con factor de escala $k_0=0.9996$ y falso este $E_0=500000$:
\[
E = k_0\,N\left(A+\frac{(1-T+C)A^3}{6}+\frac{(5-18T+T^2+72C-58e'^2)A^5}{120}\right)+E_0
\]
\[
N_{\text{UTM}} = k_0\left[M+N\tan\varphi\left(\frac{A^2}{2}+\frac{(5-T+9C+4C^2)A^4}{24}+\frac{(61-58T+T^2+600C-330e'^2)A^6}{720}\right)\right]
\]

\section{Implementación del programa}
La aplicación se implementó en Python. El programa admite dos modos de operación:
\begin{itemize}[leftmargin=1.2em]
\item \textbf{Modo puerto serie:} lectura del receptor GPS conectado al puerto \texttt{COM9} a \texttt{4800} baudios.
\item \textbf{Modo archivo:} lectura de tramas desde \texttt{prueba.txt}.
\end{itemize}

Por simplicidad (según el enunciado y el caso de uso en España), el huso UTM está fijado a \textbf{30}.

\subsection{Validación de checksum}
Cada trama NMEA incluye un checksum XOR. Para evitar procesar líneas corruptas, se valida el checksum antes de parsear campos.

\subsection{Salida}
Para cada fix con calidad distinta de 0, se imprime:
\begin{itemize}[leftmargin=1.2em]
\item Hora UTC, calidad del fix, satélites, HDOP y altitud.
\item Latitud/longitud en grados decimales.
\item Coordenadas UTM en el huso 30 para \textbf{WGS84} y \textbf{ED50}.
\end{itemize}

\section{Resultados con datos de prueba}
Se procesó el archivo \texttt{prueba.txt} (5 tramas GGA). A modo de ejemplo, para la primera trama:
\begin{itemize}[leftmargin=1.2em]
\item $\varphi = 40.386310^\circ$, $\lambda = -3.631081^\circ$ (W).
\item UTM (WGS84): $E=446435.724$~m, $N=4470826.437$~m (30N).
\item UTM (ED50): $E=446433.297$~m, $N=4470903.571$~m (30N).
\end{itemize}

La Tabla~\ref{tab:res} muestra el resumen de los 5 fixes del archivo.

\begin{table}[h]
\centering
\renewcommand{\arraystretch}{1.1}
\begin{tabular}{@{}lrr|rr@{}}
\toprule
\textbf{UTC} & \textbf{Lat (°)} & \textbf{Lon (°)} & \textbf{$E$ WGS84 (m)} & \textbf{$N$ WGS84 (m)} \\
\midrule
11:22:23.60 & 40.386310 & -3.631081 & 446435.724 & 4470826.437 \\
11:22:23.70 & 40.386310 & -3.631082 & 446435.683 & 4470826.417 \\
11:22:23.80 & 40.386310 & -3.631082 & 446435.675 & 4470826.423 \\
11:22:23.90 & 40.386310 & -3.631081 & 446435.711 & 4470826.419 \\
11:22:24.00 & 40.386310 & -3.631081 & 446435.692 & 4470826.428 \\
\bottomrule
\end{tabular}
\caption{Resultados UTM para WGS84 (huso 30N) a partir de \texttt{prueba.txt}.}
\label{tab:res}
\end{table}

\vspace{0.5em}
Como comprobación adicional, se observó un desplazamiento sistemático entre WGS84 y ED50 (del orden de decenas de metros), coherente con el cambio de datum.

\section{Cómo ejecutar}
\begin{lstlisting}[style=pythonEstilo]
# Modo archivo (lee prueba.txt)
python practica1.py --file

# Modo puerto serie (COM9, 4800 baudios)
python practica1.py --port
\end{lstlisting}

\section{Dependencias}
\begin{itemize}[leftmargin=1.2em]
\item Python 3.
\item \texttt{pyserial} (solo necesario en modo puerto): \texttt{pip install pyserial}.
\end{itemize}

\end{document}
